\section{AVALIAÇÃO DA APRENDIZAGEM}

A avaliação é um processo minucioso que permeia todo o processo ensino-aprendizagem. Avaliar não consiste somente em fazer provas e atribuir notas, mas é um processo pedagógico contínuo, que ocorre dia após dia, buscando corrigir erros e construir novos conhecimentos.

Avaliar consiste em analisar o desempenho do aluno quanto ao domínio das competências previstas face ao perfil necessário à sua formação, através da adoção de vários instrumentos e técnicas de avaliação, que deverão estar diretamente ligados ao contexto da área objeto e utilizados de acordo com a natureza do que está sendo avaliado.

A avaliação da aprendizagem acontece para que o professor tenha noção se os conhecimentos e competências necessárias à formação foram internalizados pelos discentes, bem como também serve para que o docente possa executar uma autoavaliação acerca de sua didática e metodologia de ensino, sendo possível, dessa forma, verificar se o caminho que está percorrendo deve ser revisto. Tratada dessa forma, a avaliação permite diagnosticar a situação do discente, em face da proposta pedagógica da instituição e orientar decisões quanto à condução da prática educativa. Portanto, o seu propósito é subsidiar a prática do professor, oferecendo pistas significativas para a definição e redefinição do trabalho pedagógico.

Como tal, a avaliação deverá ser contínua, processual e cumulativa, considerando a prevalência de aspectos qualitativos sobre os quantitativos, assim como estabelece a Lei de Diretrizes e Base da Educação (n° 9.394/96), isso para que seja efetivada a sua função formativa, servindo para o discente como parâmetro de referência de suas conquistas, dificuldades e possibilidades de crescimento e tendo em vista que o desenvolvimento de competências não envolve apenas conteúdos teóricos, mas, sobretudo práticas e atitudes
Nesse contexto, o processo de avaliação do \NOMECURSO do Campus de Maracanaú é orientado pelos objetivos do curso e perfil profissional do egresso a ser formado, sendo definido no Programa de Unidade Didática de cada disciplina (PUD).

Nesse sentido, vale destacar que os processos, instrumentos, critérios e valores de avaliação adotados pelo professor deverão ser explicitados aos estudantes no início do período letivo, quando da apresentação do PUD, observadas as normas dispostas no Regulamento da Organização Didática (ROD) do instituto, onde estão definidos os critérios para atribuição de notas, as formas de recuperação, promoção e frequência do aluno, assim como na Lei de Diretrizes e Base da Educação Nacional.

Considerando que o desenvolvimento de competências envolve conhecimentos, práticas e atitudes, o processo avaliativo exige diversidade de instrumentos e técnicas de avaliação, que deverão estar diretamente ligados ao contexto da área objeto da educação profissional e utilizados de acordo com a natureza do que está sendo avaliado.

Desta forma, são utilizados instrumentos diversificados que possibilitam ao professor observar e intervir no desempenho do aluno considerando os aspectos que necessitem ser melhorados, orientando a este, no percurso do curso, diante das dificuldades de aprendizagem apresentadas, reconhecendo as formas diferenciadas de aprendizagem, em seus diferentes processos, ritmos, lógicas, exercendo, assim, o seu papel de orientador e mediador que reflete na ação e que age sobre a realidade. 

Serão considerados instrumentos de avaliação, os trabalhos de natureza teórico-práticos; observação diária dos estudantes pelos professores, durante a aplicação de suas diversas atividades, exercícios, trabalhos individuais e/ou coletivos; fichas de observações; relatórios; autoavaliação; provas escritas com ou sem consulta; provas práticas e provas orais; seminários; projetos interdisciplinares, resolução de exercícios, planejamento e execução de experimentos ou projetos, relatórios referentes a trabalhos; experimentos ou visitas técnicas; realização de eventos ou atividades abertas à comunidade; autoavaliação descritiva e outros instrumentos de avaliação considerando o seu caráter progressivo e que enfatizem a resolução de situações problema específicas do processo de formação do técnico.

Dentre esses vários instrumentos podemos destacar:
\begin{itemize}
	\item Trabalho de pesquisa/projetos para verificar as capacidades de construir objetivos e alcança-los; caracterizar o que vai ser trabalhado; antecipar resultados; escolher estratégias mais adequadas à resolução do problema; executar ações; avaliar essas ações e as condições de execução; seguir critérios preestabelecidos;
	\item Observação da resolução de problemas relacionados ao trabalho em situações simuladas ou reais, com o fim de verificar que indicadores demonstram a aquisição de competências mediante os critérios de avaliação previamente estabelecidos;
	\item Seminários de exposição de conteúdos ou experiência prática de campo são procedimentos metodológicos importantes porque pressupõem o uso de ferramentas e técnicas para pesquisa, estudo e trabalho em equipe;
	\item Análise de casos – os casos são desencadeadores de um processo de pensar, fomentador da dúvida, do levantamento e da comprovação de hipóteses, do pensamento inferencial, do pensamento divergente, entre outros.
	\item Prova – visa verificar a capacidade adquirida pelos alunos de aplicar os conteúdos aprendidos. Como, por exemplo: analisar, classificar, comparar, criticar, generalizar e levantar hipóteses, estabelecer relações com base em fatos, fenômenos, ideias e conceitos. Para fins de promoção são avaliados tanto o desempenho quanto a assiduidade do aluno. O aluno será orientado na medida em que os resultados das atividades não sejam apenas comunicados, mas discutidos, indicando erros, identificando dificuldades e limitações, sugerindo possíveis soluções e rumos, considerando o caráter progressivo da avaliação.
	A sistemática de avaliação no IFCE se desenvolverá em duas etapas e em cada uma será computada a média obtida pelo discente. Independentemente do número de aulas semanais, o docente deverá aplicar, no mínimo, duas avaliações por etapa. 
\end{itemize}

O estudante que não atingir o mínimo necessário para aprovação, poderá realizar avaliação de recuperação, conforme estabelecido no Regulamento da Organização Didática do IFCE. A sistemática de avaliação no IFCE é apresentada na Subseção I, Seção I, Capítulo III, Título III, do Regulamento da Organização Didática (ROD) de junho de 2015, ANEXO \ref{anexoRODavaliacao}.
